% Problem record op_94fc1874fe23df16. This TeX file is derived from
% database/problems_json/op_94fc1874fe23df16.json by scripts/sync-tex.mjs; edit the JSON record
% and rerun the script rather than editing this file.
\section{Average-case approximation hardness of random-circuit output probabilities}
\label{sec:op_94fc1874fe23df16}

\paragraph{Problem.}
Does there exist a fixed family of $n$-qubit circuit layouts with
$m=\operatorname{poly}(n)$ one- and two-qubit gates for which the following
task is $\#\mathrm{P}$-hard? Choose every gate independently from Haar
measure, obtaining $C$. Given $C$ and $\epsilon,\delta\in(0,1)$, estimate
the all-zero output probability with
\begin{equation}
  p_0(C)=|\langle0^n|C|0^n\rangle|^2,\qquad
  \Pr_C[|\widetilde p(C)-p_0(C)|\leq\epsilon 2^{-n}]\geq1-\delta.
  \label{eq:rcs-average-approximation}
\end{equation}
The guarantee in Eq.~\eqref{eq:rcs-average-approximation} must hold for
arbitrary $\epsilon,\delta$, with running time measured in
$n,1/\epsilon,1/\delta$.

\subsection*{Status}

Unsolved

\subsection*{Source}

Conjecture 6 of Bouland, Fefferman, Nirkhe, and Vazirani (2018), in its
formal version in Appendix A.4, using Definition 20 of an average-case
approximate solution \sourcecite{ref:rcs-bfnv}{BFNV19}.
This is an additive-error probability-estimation conjecture; the parameters
$\epsilon,\delta$ retain the original meaning.

\subsection*{Progress}

\begin{itemize}
  \item 2018: Taylor truncation gives exact $\#\mathrm{P}$-hardness on an $8/9$
fraction of perturbed, generally nonunitary instances (Theorem 1);
it does not establish the desired approximation guarantee for Haar gates
\sourcecite{ref:rcs-bfnv}{BFNV19}.

  \item 2019--2023: Movassagh's Cayley path preserves unitarity and enables rational
interpolation, proving exact $\#\mathrm{P}$-hardness on a
$3/4+1/\operatorname{poly}(n)$ fraction of Haar-random circuits
(arXiv Theorem 1) \sourcecite{ref:rcs-movassagh}{Mov23}.

  \item 2021: Bouland et al. and, independently, Kondo--Mori--Movassagh obtain
additive-error tolerance $2^{-O(m\log m)}$ with constant failure probability,
using $\mathrm{BPP}^{\mathrm{NP}}$ reductions
\sourcecite{ref:rcs-bfll}{BFLL22}\sourcecite{ref:rcs-kmm}{KMM22}.

  \item 2022: Krovi improves the tolerance to $2^{-O(m)}$ for a constant fraction
of Haar-random circuits, via $\mathrm{BPP}$ reductions; the hardness class
is $\mathrm{coC}_{=}\mathrm{P}$, rather than $\#\mathrm{P}$
\sourcecite{ref:rcs-krovi}{Kro22}.

  \item 2025: Bouland et al.'s dilution method gives $\#\mathrm{P}$-hardness at
error $2^{-n-O(n^\gamma)}$ for suitable depth-$\Omega(\log n)$ RCS ensembles,
for every fixed $\gamma>0$ (Corollary 2, $\mathrm{BPP}^{\mathrm{NP}}$ reduction)
\sourcecite{ref:rcs-bdfh}{BDFH25}.
\end{itemize}

\subsection*{References}

\begin{enumerate}[label={},leftmargin=4.8em,labelsep=0.5em]
  \footnotesize
  \item[\textup{[BFNV19]}]\label{ref:rcs-bfnv}
  A. Bouland, B. Fefferman, C. Nirkhe, and U. Vazirani,
"On the complexity and verification of quantum random circuit sampling,"
\emph{Nature Physics} \textbf{15}, 159--163 (2019).
\href{https://doi.org/10.1038/s41567-018-0318-2}{doi:10.1038/s41567-018-0318-2};
\href{https://arxiv.org/abs/1803.04402}{arXiv:1803.04402} (2018 preprint titled
"Quantum Supremacy and the Complexity of Random Circuit Sampling").

  \item[\textup{[Mov23]}]\label{ref:rcs-movassagh}
  R. Movassagh, "The hardness of random quantum circuits,"
\emph{Nature Physics} \textbf{19}, 1719--1724 (2023).
\href{https://doi.org/10.1038/s41567-023-02131-2}{doi:10.1038/s41567-023-02131-2};
\href{https://arxiv.org/abs/1909.06210v4}{arXiv:1909.06210v4}
(preprint titled "Quantum supremacy and random circuits").

  \item[\textup{[BFLL22]}]\label{ref:rcs-bfll}
  A. Bouland, B. Fefferman, Z. Landau, and Y. Liu,
"Noise and the frontier of quantum supremacy,"
in \emph{FOCS 2021}, 1308--1317 (2022).
\href{https://doi.org/10.1109/FOCS52979.2021.00127}{doi:10.1109/FOCS52979.2021.00127};
\href{https://arxiv.org/abs/2102.01738v2}{arXiv:2102.01738v2}.

  \item[\textup{[KMM22]}]\label{ref:rcs-kmm}
  Y. Kondo, R. Mori, and R. Movassagh,
"Quantum supremacy and hardness of estimating output probabilities of quantum circuits,"
in \emph{FOCS 2021}, 1296--1307 (2022).
\href{https://doi.org/10.1109/FOCS52979.2021.00126}{doi:10.1109/FOCS52979.2021.00126};
\href{https://arxiv.org/abs/2102.01960v3}{arXiv:2102.01960v3}.

  \item[\textup{[Kro22]}]\label{ref:rcs-krovi}
  H. Krovi, "Average-case hardness of estimating probabilities of random
quantum circuits with a linear scaling in the error exponent,"
arXiv preprint (2022).
\href{https://doi.org/10.48550/arXiv.2206.05642}{doi:10.48550/arXiv.2206.05642};
\href{https://arxiv.org/abs/2206.05642}{arXiv:2206.05642}.

  \item[\textup{[BDFH25]}]\label{ref:rcs-bdfh}
  A. Bouland, I. Datta, B. Fefferman, and F. Hernandez,
"Exponential improvements to the average-case hardness of BosonSampling,"
in \emph{FOCS 2025}, 912--933 (2025).
\href{https://doi.org/10.1109/FOCS63196.2025.00047}{doi:10.1109/FOCS63196.2025.00047};
\href{https://arxiv.org/abs/2411.04566v2}{arXiv:2411.04566v2}.
\end{enumerate}

\subsection*{Comment}

The original approximation conjecture remains open. The 2025 result still
has an $O(n^\gamma)$ loss in the exponent, rather than the $O(\log n)$
loss corresponding to error $2^{-n}/\operatorname{poly}(n)$
\sourcecite{ref:rcs-bdfh}{BDFH25}. Exact hardness and smaller additive-error
hardness do not settle this gap.

\subsection*{Field}

Quantum algorithm

\subsection*{Topic}

Computational complexity; Quantum supremacy; Random circuit sampling

\subsection*{ID}

\texttt{op\_94fc1874fe23df16}
